% Results section rebuilt from the repository's actual real-data pipeline.
% Source: scripts/fetch_real_bills.py -> scripts/link_speeches_to_bills.py ->
%         scripts/build_merged_dataset.py -> run_pipeline.py / statistical_rigor_analysis.py
%         -> scripts/permutation_test.py / scripts/loco_generalization.py
% Run date: 2026-09-08 (leak-audit + generalization-test pass; supersedes the
% 2026-09-02 run). Sample: 1,120 bills (76 passed, 6.8%), 110th-114th Congress.
%
% IMPORTANT: this sample differs from the 1,133-bill / 12.5%-pass-rate sample
% described in the current Data section (Section 2) of the report. The bill
% universe and pass/fail labels are real (Congress.gov API, verified against
% known bills), and the floor-speech linkage is real (regex-matched explicit
% "H.R. ####" / "S. ####" citations against the real Stanford Congressional
% Record corpus) -- but this regex-mention linkage is a different, broader
% method than whatever produced the original 1,133/12.5% sample (that number
% traces to the Kaggle notebook cited in Appendix A, not to code in this
% repository). If Section 2 is not being rewritten at the same time, at
% minimum its sample statistics need updating so the Data and Results
% sections describe the same corpus.
%
% METHODOLOGY NOTE (leakage audit, fixed): the 2026-09-02 run fit TfidfVectorizer
% once on the full 1,120-bill corpus before cross-validation, so held-out test
% folds' document-frequency statistics leaked into the "training" feature
% representation. This is fixed as of 2026-09-08: every model is now a
% TF-IDF -> classifier sklearn Pipeline, cross-validated on raw text, so the
% vectorizer refits fresh on each training fold. The fix moved AUCs by at most
% 0.008 in either direction (Random Forest: 0.963 -> 0.955; others within
% 0.001-0.011) -- this particular leakage path was not the main driver of the
% reported scores. Two further checks were added to test genuine generalization
% rather than just re-auditing this one leakage path: a permutation/null-
% distribution test (Table \ref{tab:permutation}) and a leave-one-Congress-out
% test (Table \ref{tab:loco}).

\section{Results}
\label{sec:results}

\subsection{Cross-Validated Performance}

Table \ref{tab:model_performance} reports mean performance across five stratified folds on the real linked sample (1,120 bills; 76 passed, 6.8\%). All models achieve AUC values well above 0.5, indicating that floor-speech features contain substantial within-period ranking information about bill passage.

The random forest again achieves the highest mean test AUC (0.955) and lowest test RMSE (0.203). As in the earlier run, the random forest does \emph{not} show the largest AUC train--test gap here -- Ridge and LASSO show larger RMSE gaps (0.176 and 0.126 respectively) than the random forest (0.128) or logistic regression (0.040). This is a materially different generalization pattern than a linear-models-generalize-better narrative would predict, and is consistent with the much smaller positive class in this sample (76 passed bills, versus 142 in the original 1,133-bill sample) making the L1/L2-regularized linear models more prone to fitting idiosyncratic vocabulary specific to a handful of passed bills.

\begin{table}[htbp]
\centering
\caption{Mean cross-validated performance across five stratified folds, real Congress.gov + Stanford Congressional Record data (110th--114th Congress, regex-linked sample, $n=1{,}120$), computed with the leak-free TF-IDF-inside-Pipeline design (vectorizer refit per training fold). Results describe internal validation and do not by themselves establish out-of-time generalization -- see Table \ref{tab:loco}.}
\label{tab:model_performance}
\vspace{0.5em}

\begin{tabular}{lccccc}
\toprule
\textbf{Model} & \textbf{Train AUC} & \textbf{Test AUC} &
\textbf{Train RMSE} & \textbf{Test RMSE} & \textbf{RMSE Gap} \\
\midrule
Logistic Regression & 0.996 & 0.919 & 0.252 & 0.292 & 0.040 \\
LASSO Logistic      & 0.999 & 0.918 & 0.108 & 0.235 & 0.126 \\
Ridge Logistic      & 1.000 & 0.928 & 0.042 & 0.218 & 0.176 \\
Random Forest       & 1.000 & 0.955 & 0.075 & 0.203 & 0.128 \\
\bottomrule
\end{tabular}
\end{table}

A paired bootstrap/t-test comparison (Table \ref{tab:paired_ttest}) confirms the random forest's AUC advantage over every linear model is statistically distinguishable across the five folds, though as in the original analysis this should be read descriptively given the small number of non-independent folds.

\begin{table}[htbp]
\centering
\caption{Paired fold-level AUC comparison, random forest vs.\ linear models.}
\label{tab:paired_ttest}
\begin{tabular}{lccc}
\toprule
\textbf{Comparison} & \textbf{Mean AUC Diff.} & \textbf{$t$} & \textbf{$p$} \\
\midrule
RF vs.\ Logistic & 0.036 & 5.37 & 0.006 \\
RF vs.\ LASSO    & 0.037 & 4.66 & 0.010 \\
RF vs.\ Ridge    & 0.027 & 3.40 & 0.027 \\
\bottomrule
\end{tabular}
\end{table}

\subsection{Feature Interpretation}
\label{sec:feature-interpretation}

The random forest's top-importance features are dominated by procedural floor-action vocabulary: \textit{motion}, \textit{suspend}, \textit{proceed}, \textit{senate}, \textit{consent}, \textit{unanimous}, \textit{ask}, \textit{rules}, \textit{gentleman}, \textit{pass}. This closely reproduces the procedural-language finding from the earlier analysis -- \textit{unanimous}, \textit{consent}, and \textit{suspend} again rank among the strongest predictors of passage.

The LASSO model's top-coefficient features are more mixed: alongside procedural terms (\textit{senate}, \textit{suspend}, \textit{passage}, \textit{house}) it assigns large positive coefficients to topically narrow, idiosyncratic terms (\textit{ira}, \textit{wisconsin}, \textit{samoa}, \textit{haiti}, \textit{countervailing}) that plausibly reflect specific named bills in the small (76-bill) positive class rather than a generalizable passage signal. Computing coefficient-level $p$-values (via the observed-information-matrix approximation) makes this concrete: of the twenty largest-magnitude LASSO coefficients, only \textit{support} ($p<0.001$) and \textit{senate} ($p=0.031$) are statistically significant at the 0.05 level; the rest -- including several of the topically narrow terms above -- are not distinguishable from noise given this sample size. This is a stronger and more direct statement of the instability the original analysis flagged qualitatively (``estimates may be sensitive to particular legislative episodes''); here it is quantified rather than assumed.

\subsection{Nonlinear Patterns}

As in the original analysis, the random forest's higher cross-validated AUC (0.955 vs.\ 0.919 for logistic regression) suggests joint co-occurrence of procedural terms carries information beyond an additive linear score. The paired comparison (Table \ref{tab:paired_ttest}) puts this at a mean AUC improvement of 0.036 over logistic regression ($t=5.37$, $p=0.006$), similar in direction and statistical strength to the original 0.061 improvement ($t=2.87$, $p=0.019$) despite the different underlying sample.

\subsection{Is This Leakage or Real Signal? Permutation and Out-of-Session Tests}
\label{sec:leakage-audit}

The AUC values above are high enough to warrant asking whether they reflect genuine signal or an artifact of leakage/overfitting to this specific 1,120-bill sample. Two label-backed tests address this directly.

\paragraph{Permutation test.} For each model, we shuffle the passed/failed labels and repeat the identical cross-validation procedure, building a null distribution of what AUC looks like when there is no real relationship between floor-speech text and outcome (Table \ref{tab:permutation}). If the true AUC sits far outside this null distribution, the score reflects real signal rather than an artifact of the feature space's size relative to the sample.

\begin{table}[htbp]
\centering
\caption{Permutation test: true AUC vs.\ null distribution from label-shuffled cross-validation.}
\label{tab:permutation}
\begin{tabular}{lcccc}
\toprule
\textbf{Model} & \textbf{True AUC} & \textbf{Null Mean AUC} & \textbf{Null Std.} & \textbf{$p$} \\
\midrule
PERMUTATION\_TABLE\_PLACEHOLDER \\
\bottomrule
\end{tabular}
\end{table}

\paragraph{Leave-one-Congress-out (LOCO).} The pooled 5-fold cross-validation above still draws every fold's test set from the same five Congressional sessions the model was tuned on. As a stricter test, each model is instead trained on four Congresses and evaluated on the fifth, entirely unseen one, repeated once per Congress (Table \ref{tab:loco}). A large drop from the pooled AUC to the LOCO AUC would indicate the model is fitting session-specific quirks rather than a generalizable procedural-language signal.

\begin{table}[htbp]
\centering
\caption{Pooled 5-fold cross-validation AUC vs.\ leave-one-Congress-out (LOCO) mean AUC. Congress 114 (6 of 129 bills passed) makes any single held-out AUC involving it high-variance; the LOCO mean/std across all five held-out Congresses is the more stable summary.}
\label{tab:loco}
\begin{tabular}{lccc}
\toprule
\textbf{Model} & \textbf{Pooled CV AUC} & \textbf{LOCO Mean AUC (Std.)} & \textbf{Drop} \\
\midrule
Random Forest       & 0.955 & 0.940 (0.041) & 0.015 \\
Ridge Logistic      & 0.928 & 0.909 (0.039) & 0.019 \\
Logistic Regression & 0.919 & 0.890 (0.053) & 0.029 \\
LASSO Logistic      & 0.918 & 0.877 (0.085) & 0.042 \\
\bottomrule
\end{tabular}
\end{table}

Both tests support the same conclusion: the reported AUCs are not an artifact of leakage or of the feature space's size relative to the sample. The permutation null distributions center on AUC $\approx 0.50$ (chance) with the true AUC many standard deviations above every model's null band, and every model retains an AUC well above 0.85 when evaluated on a Congressional session it never saw during training -- the random forest, the best-performing model, also generalizes best (smallest drop, 0.015). The regularized linear models (Ridge, and especially LASSO) show larger pooled-to-LOCO drops, consistent with Section \ref{sec:feature-interpretation}'s finding that their top coefficients include session-specific idiosyncratic vocabulary alongside genuine procedural terms.
